% PROPOSED, unlabeled. Becomes notes/E##-model-ceiling.tex when Enrico agrees. Compile: latexmk -pdf
\documentclass[11pt]{article}
\usepackage{enrico}
\usepackage{booktabs}

\title{Is the ceiling the model's or the game's?}
\author{Enrico Yao-Bate}
\date{design written 2026-10-08; not yet agreed; no run started}

\begin{document}
\maketitle

\section{Question}
The local model cannot improve the 17-point reference bot: 0 of 30 single revisions, and two 40-generation
chains never left its behaviour (\texttt{findings/2026-10-07-local-model-ceiling.md}). Every accumulation
experiment is bounded above by the innovation rate at the level the population has reached, so before any
population run on a new model we need its rate at three levels. This is a measurement, not a hypothesis test,
but it has a decision attached: whether to run populations on this model at all, and whether the revise prompt
is the bottleneck.

\section{Setting}
Three incumbents: the weak student (held-out $2.75$), the 16-point bot (IGGI, $15.8$), the 17-point bot
(Piers, $17.0$). One revision each, nothing delivered, feedback traces varying across 30 replicates (experiment
seeds $400$--$429$, as in the local probe), verification off and no adoption, held-out paired evaluation on 300
deals and a common 300-deal set so levels are comparable. Hosted model, temperature 0, seed fixed (providers may
not honour it; replicates are seeds, not replays).

\section{What we measure}
For incumbent level $\ell$, the innovation rate $\epsilon(\ell)=\P(\Delta>0.5)$ with $\Delta$ the paired
held-out improvement, Beta(1,1) posterior and 95\% interval; the gain distribution given a gain; the no-op
rate (candidate byte-identical to the parent); the inadmissible rate (fails the sandbox or the smoke game).
The no-op rate is the diagnostic for ``the prompt tells the model to leave good code alone'' versus ``the model
cannot find an improvement''.

\section{Prediction (2026-10-08)}
For a flash-class open-weight model: $\epsilon(2.75)\ge 0.6$ with gains around $8$; $\epsilon(15.8)\in[0.1,0.3]$;
$\epsilon(17.0)\in[0.05,0.2]$ with gains of $0.5$ to $2$; no-op rate below $0.3$. If $\epsilon(17.0)$ is again
$0/30$ and the no-op rate is above $0.5$, the ceiling is in the prompt (it asks for a rewrite of the worst games'
causes and the 17-point bot has few), and the next step is a prompt variant, not a bigger model.

\section{Design}
Three levels $\times$ 30 replicates, 90 calls, paired on experiment seeds with the local probe so the two
models are compared on the same traces. Decision rule, fixed now: populations run on this model only if the
lower 95\% bound of $\epsilon(17.0)$ exceeds $0.03$; otherwise the prompt is revised and the probe re-run before
any population spend. Figure: $\epsilon$ against level for the two models, with intervals; gain histograms.

\section{Cost}
90 calls, about one dollar at flash-class prices, twenty minutes in parallel. This is also the first paid
probe of the provider: it reports price per call, latency and whether the provider is deterministic at
temperature 0 with a seed.

\section{Nearest prior work}
Innovation-rate ladders for program-writing agents were not found; the closest are self-improvement
curves in evolutionary code search (AlphaEvolve-style systems report best-so-far against iterations, not a rate
by level). Phrase accordingly.
\end{document}
